% =====================================================================
%  FICHE 4 — THEME 1 — Corrige niveau DIFFICILE
% =====================================================================
\documentclass[11pt,a4paper]{article}
\usepackage[utf8]{inputenc}
\usepackage[T1]{fontenc}
\usepackage{lmodern}
\usepackage[french]{babel}
\usepackage{amsmath,amssymb,mathtools}
\usepackage{icomma}
\usepackage[version=4]{mhchem}
\usepackage{siunitx}
\sisetup{output-decimal-marker={,},per-mode=symbol}
\DeclareSIUnit\Molar{M}
\usepackage{xcolor}
\usepackage{tikz}
\usepackage[most]{tcolorbox}
\usepackage{enumitem}
\usepackage{booktabs}
\usepackage{array}
\usepackage{fancyhdr}
\usepackage[a4paper,margin=2.2cm,headheight=26pt,headsep=10pt]{geometry}

\definecolor{primary}{RGB}{146,95,160}
\definecolor{primarydark}{RGB}{92,52,104}
\definecolor{excol}{RGB}{0,135,90}

\tcbset{boxbase/.style={enhanced,breakable,boxrule=0.9pt,arc=2.5pt,
   left=3mm,right=3mm,top=2mm,bottom=2mm,fonttitle=\bfseries,coltitle=white,
   toptitle=0.5mm,bottomtitle=0.5mm}}

\newcounter{exo}
\newenvironment{corr}[1][]{%
  \par\medskip\refstepcounter{exo}%
  \noindent\begin{tcolorbox}[boxbase,colback=excol!5,colframe=excol,
     title={\textsc{Correction \theexo}\ifx\relax#1\relax\relax\else\textmd{ --- \itshape #1}\fi}]}%
  {\end{tcolorbox}}

\pagestyle{fancy}\fancyhf{}
\fancyhead[L]{\small\textcolor{primarydark}{\textbf{Fiche 4} --- Th.\,1 : Mole, masse molaire, entit\'es}}
\fancyhead[R]{\small\textcolor{primarydark}{\textsc{Chimie} --- \textbf{Corrigé Difficile}}}
\fancyfoot[C]{\small\thepage}
\renewcommand{\headrulewidth}{0.8pt}
\renewcommand{\headrule}{\hbox to\headwidth{\color{primary}\leaders\hrule height \headrulewidth\hfill}}
\newcommand{\NA}{N_{\!\text{A}}}
\setlength{\parindent}{0pt}

\begin{document}

\begin{center}
{\LARGE\bfseries\color{primarydark} Thème 1 --- Corrigé Difficile}
\end{center}

\begin{corr}[eau de cristallisation d'un hydrate]
\textbf{a)} $M(\ce{NiCl2})=58{,}7+2(35{,}5)=\SI{129.7}{\gram\per\mole}$. Le solide restant est le
\ce{NiCl2} anhydre :
\[ n(\ce{NiCl2})=\frac{0{,}130}{129{,}7}=\SI{1.00e-3}{\mole}. \]
\textbf{b)} L'eau partie est la différence de masse (hydrate $-$ anhydre) :
\[ m(\ce{H2O})=0{,}238-0{,}130=\SI{0.108}{\gram}\ \Rightarrow\
   n(\ce{H2O})=\frac{0{,}108}{18}=\SI{6.0e-3}{\mole}. \]
\textbf{c)} $x$ est le nombre de moles d'eau \emph{par} mole de sel :
\[ x=\frac{n(\ce{H2O})}{n(\ce{NiCl2})}=\frac{\num{6.0e-3}}{\num{1.00e-3}}=6. \]
\[ \boxed{x=6\ \Rightarrow\ \ce{NiCl2.6H2O}} \]
\textit{Idée-clé :} l'eau perdue au chauffage est exactement l'eau de cristallisation ; le rapport
des quantités de matière donne directement $x$.
\end{corr}

\begin{corr}[identifier un métal, puis nommer son chlorure]
\textbf{a)} La formule \ce{M2O} contient deux atomes de métal et un oxygène :
\[ M(\ce{M2O})=2\,M(\ce{M})+M(\ce{O})\ \Rightarrow\ 2\,M(\ce{M})+16=232. \]
\[ 2\,M(\ce{M})=216\ \Rightarrow\ M(\ce{M})=\SI{108}{\gram\per\mole}\ \Rightarrow\ \text{c'est l'argent \ce{Ag}}. \]
\textbf{b)} Le chlorure d'argent \ce{MCl}$=$\ce{AgCl} :
\[ M(\ce{AgCl})=108+35{,}5=\SI{143.5}{\gram\per\mole}. \]
\[ \boxed{M=\ce{Ag};\quad M(\ce{AgCl})=\SI{143.5}{\gram\per\mole}} \]
\end{corr}

\begin{corr}[pourcentage massique utilisé « à l'envers »]
\textbf{a)} Dans une mole de chlorophylle il y a une mole d'atomes de \ce{Mg}, soit $\SI{24}{\gram}$.
Cette masse représente $\SI{3.0}{\percent}$ de la masse molaire $M$ :
\[ \%\,\ce{Mg}=\frac{M(\ce{Mg})}{M}\times100\ \Rightarrow\ 3{,}0=\frac{24}{M}\times100
   \ \Rightarrow\ M=\frac{24}{0{,}030}=\SI{800}{\gram\per\mole}. \]
\textbf{b)} $n=\dfrac{1{,}0}{800}=\SI{1.25e-3}{\mole}$, d'où
\[ N=n\times\NA=\num{1.25e-3}\times\num{6.02e23}\approx\num{7.5e20}\ \text{molécules}. \]
\[ \boxed{M\approx\SI{800}{\gram\per\mole};\quad N\approx\num{7.5e20}\ \text{molécules}} \]
\textit{Idée-clé :} le pourcentage massique relie une masse \emph{connue} (celle du \ce{Mg}) à la
masse molaire \emph{cherchée}.
\end{corr}

\begin{corr}[identifier un élément alcalin]
\textbf{a)} La masse molaire se déduit de $n=m/M$ :
\[ M(\ce{M3PO4})=\frac{m}{n}=\frac{42{,}4}{0{,}20}=\SI{212}{\gram\per\mole}. \]
\textbf{b)} La formule \ce{M3PO4} contient trois atomes de métal, un phosphore et quatre oxygènes :
\[ 3\,M(\ce{M})+31+4(16)=212\ \Rightarrow\ 3\,M(\ce{M})+95=212\ \Rightarrow\ 3\,M(\ce{M})=117. \]
\[ M(\ce{M})=\frac{117}{3}=\SI{39}{\gram\per\mole}\ \Rightarrow\ \text{c'est le potassium \ce{K}}. \]
\[ \boxed{\ce{M}=\ce{K}\ (\text{potassium}),\ \text{le composé est }\ce{K3PO4}} \]
\textit{Méthode générale :} passer par la masse molaire du composé, puis isoler $M(\ce{M})$ dans la
formule, enfin comparer aux valeurs des alcalins.
\end{corr}

\end{document}
